\chapter[1983 --- McClintock]{1983: Barbara McClintock}
\label{ch:1983_barbaramcclintock}

\section*{Main Contribution}

\textit{Discovered transposable genetic elements (jumping genes), showing that genes can move within the genome---challenging the view that genes are fixed in place.}

\section{Official Citation}

for her discovery of mobile genetic elements

The 1983 Nobel Prize in Physiology or Medicine was awarded to Barbara McClintock ``for her discovery of mobile genetic elements.'' McClintock's discovery of transposable elements---genetic elements that can move from one position to another within the genome---was one of the most radical and far-reaching insights in the history of genetics. It challenged the prevailing view that the genome was a stable, fixed structure and revealed that genes could change their position, alter their expression, and influence the phenotype of the organism in unexpected and dynamic ways.

\section{Background and Historical Context}

Barbara McClintock was born in Hartford, Connecticut, in 1902, and grew up in a family that encouraged her intellectual curiosity and independence. She attended Cornell University, where she earned her bachelor's degree in 1923, her master's degree in 1925, and her Ph.D. in botany in 1927. At Cornell, McClintock studied under Rollins A. Emerson, a pioneering geneticist who was studying the genetics of maize (\emph{Zea mays})---a model organism that had been used in genetic studies since the work of Walter Sutton and Edward M. East in the early twentieth century.

McClintock's early work at Cornell and later at the University of Missouri and Cold Spring Harbor Laboratory established her as one of the leading cytogeneticists of her generation. She developed techniques for visualizing and identifying individual chromosomes in maize cells, and she used these techniques to correlate specific chromosomal features with specific genetic traits. In the late 1930s and early 1940s, McClintock made several important discoveries about the structure and behavior of chromosomes, including the identification of the nucleolar organizer region, the discovery of chromosomal crossover (recombination) during meiosis, and the demonstration that broken chromosomes could fuse to form ring chromosomes.

The historical context of McClintock's discovery of transposable elements was shaped by the prevailing orthodoxy of classical genetics. In the 1940s and 1950s, the dominant view was that genes were fixed entities located at specific positions (loci) on chromosomes, and that the arrangement of genes on chromosomes was stable and unchanging. The idea that genes could move from one position to another within the genome was considered heretical---a violation of the fundamental principles of genetics as they were then understood.

McClintock's interest in transposable elements began with her study of maize kernels, which display a remarkable variety of colors and patterns. The pigmentation of maize kernels is determined by the expression of genes involved in the biosynthesis of anthocyanin pigments, and these genes are subject to a variety of regulatory mechanisms that control their expression in different tissues and at different times during kernel development. McClintock observed that certain maize strains displayed unusual pigmentation patterns---such as kernels with sectors of color colorless tissue---that suggested that the genes controlling pigmentation were subject to sudden, heritable changes in expression.

\section{The Discovery/Contribution}

\subsection{The Discovery of the Ac/Ds System}

In the late 1940s, McClintock conducted a series of elegant genetic experiments that led to her discovery of the first transposable elements in maize. Using a combination of genetic crosses, cytological analysis, and careful observation of kernel phenotypes, McClintock identified two genetic elements that she named Activator (Ac) and Dissociation (Ds). She showed that Ds was a genetic element that could cause breakage of the chromosome at a specific locus, and that Ac was a regulatory element that could activate (or ``mobiliate'') Ds, causing it to move from one chromosomal position to another.

The key observation was this: when Ds was located near a gene involved in pigment biosynthesis, the presence of Ac caused the gene to be inactivated in some cells and active in others, producing a mosaic or variegated pattern of kernel pigmentation. McClintock realized that this variegation was caused by the movement of Ds into and out of the pigment gene. When Ds inserted into the gene, the gene was inactivated; when Ds excised (left) the gene, the gene was reactivated. The timing of these insertions and excisions during kernel development determined the size and pattern of the pigmented sectors.

McClintock published her first report on the Ac/Ds system in 1948, and she continued to develop and refine her ideas over the following decades. She showed that transposable elements could move to new chromosomal positions, that they could alter the expression of nearby genes, and that their movement was regulated by complex genetic interactions between the transposable elements themselves and the host genome.

\subsection{The Spm/En System and Beyond}

In the 1950s, McClintock discovered a second transposable element system in maize, which she called the Spm (Suppressor-mutator) element. The Spm system was more complex than the Ac/Ds system, and it displayed a wider range of effects on gene expression. McClintock showed that Spm could not only activate and inactivate genes but could also modulate their expression levels in a graded, dose-dependent manner. She described Spm as having both ``regulatory'' and ``mutable'' functions---it could regulate the expression of genes at a distance and, at the same time, it could itself be subject to changes in activity.

McClintock's work on transposable elements was based entirely on genetic and cytological analysis---she did not have access to the molecular tools that would later be used to characterize transposable elements at the DNA level. Her methods were painstaking and labor-intensive: she grew thousands of maize plants, crossed them in specific combinations, and analyzed the resulting kernel phenotypes under the microscope. Her ability to detect subtle patterns of variation and to interpret them in terms of underlying genetic mechanisms was notable, and it reflected her deep knowledge of maize genetics and her exceptional powers of observation.

\subsection{The Struggle for Recognition}

McClintock's discovery of transposable elements was met with skepticism and incomprehension by much of the genetics community. Her first report on the Ac/Ds system was presented at a Cold Spring Harbor Symposium in 1948, and while it attracted some interest, it was not widely understood or accepted. Many geneticists regarded transposable elements as a peculiarity of maize that was unlikely to be relevant to other organisms. The prevailing view was that the genome was a stable, fixed structure, and the idea that genes could move around within it seemed to contradict the fundamental principles of Mendelian genetics.

McClintock's frustration with the lack of recognition for her work is well documented. In interviews and autobiographical writings, she described feeling isolated and misunderstood, and she often expressed bewilderment at the inability of other scientists to appreciate the significance of her discoveries. Despite the lack of recognition, McClintock continued her work with unwavering dedication, publishing a comprehensive review of her transposition experiments in \emph{Cold Spring Harbor Symposia on Quantitative Biology} in 1951 and continuing to present her findings at scientific meetings and in published papers throughout the 1950s and 1960s.

The turning point came in the 1970s, when molecular biologists discovered transposable elements in bacteria (the insertion sequences and transposons of \emph{E. coli}), yeast, \emph{Drosophila}, and other organisms. These molecular discoveries validated McClintock's genetic observations and demonstrated that transposable elements were a universal feature of genomes, not a peculiarity of maize. The molecular characterization of transposable elements revealed that they are DNA sequences that encode the enzymes necessary for their own transposition, and that they can move from one chromosomal position to another by a variety of mechanisms, including ``cut and paste'' transposition (in which the element is excised from one position and inserted into another) and ``copy and paste'' transposition (in which the element is replicated, and the copy is inserted into a new position).

By the early 1980s, it had become clear that McClintock's discoveries were of fundamental importance for genetics, and in 1983, she was awarded the Nobel Prize in Physiology or Medicine---the first woman to receive the prize without sharing it. She was 81 years old at the time, and she continued to work at Cold Spring Harbor Laboratory until shortly before her death in 1992.

\section{Impact on Modern Medicine}

The impact of McClintock's discovery of transposable elements on modern medicine has been enormous and multifaceted. Transposable elements have been found to play important roles in a wide range of biological processes and diseases, and they have become powerful tools for genetic research and gene therapy.

In genetics, transposable elements have been recognized as major drivers of genome evolution. Approximately 45\% of the human genome consists of sequences derived from transposable elements, including long interspersed nuclear elements (LINEs), short interspersed nuclear elements (SINEs), long terminal repeat (LTR) retrotransposons, and DNA transposons. These elements have shaped the structure and function of the human genome over millions of years, contributing to gene duplication, gene regulation, and the creation of new genes.

Transposable elements have also been implicated in a variety of human diseases. Insertional mutagenesis---the disruption of a gene by the insertion of a transposable element---has been shown to cause genetic diseases such as hemophilia A (caused by the insertion of an LINE-1 element into the factor VIII gene), neurofibromatosis type 1, and certain forms of muscular dystrophy. Somatic transposition---the movement of transposable elements in non-reproductive tissues---has been implicated in cancer, and there is growing evidence that transposable elements contribute to the genomic instability that is a hallmark of many tumors.

In gene therapy, transposable elements have been adapted as vectors for the delivery of therapeutic genes. Sleeping Beauty, a DNA transposon that was reconstructed from inactive transposon sequences found in fish genomes, has been used in clinical trials for the treatment of cancer and genetic diseases. piggyBac, another DNA transposon, has also been developed as a gene therapy vector. These transposon-based vectors offer several advantages over viral vectors, including greater safety (because they do not carry viral sequences that could cause immune reactions) and the ability to integrate genes into the host genome in a more predictable manner.

Transposable elements have also been exploited as tools for genetic research. The \emph{P} element of \emph{Drosophila melanogaster} has been used as a mutagen and gene-tagging vector in fruit fly genetics for decades, and transposon-based mutagenesis screens have been used to identify genes involved in development, behavior, and disease in a wide range of organisms.

More recently, transposable elements have been recognized as important regulators of gene expression. The insertion of transposable elements near genes can alter their expression by providing new promoter or enhancer sequences, and there is evidence that some transposable elements have been ``co-opted'' by the host genome to serve as regulatory elements. For example, approximately 25\% of the promoter regions of human genes contain sequences derived from transposable elements, suggesting that transposable elements have played a major role in the evolution of gene regulation.

In cancer biology, transposable elements have been implicated in several aspects of tumor development and progression. The expression of transposable elements is often derepressed in cancer cells, leading to increased transposition activity and genomic instability. This derepression has been linked to the silencing of tumor suppressor genes, the activation of oncogenes, and the generation of chromosomal rearrangements that are characteristic of many cancers. There is also evidence that the immune system can recognize and respond to transposable element-derived sequences expressed in cancer cells, raising the possibility that transposable elements could serve as targets for cancer immunotherapy.

\section{Legacy}

Barbara McClintock's legacy is one of a notable in the history of science. Her discovery of transposable elements was a triumph of independent thinking and meticulous observation---a discovery made by a single scientist working largely alone, against the prevailing orthodoxy, using methods that were considered old-fashioned by the molecular biologists of her day. The fact that her discovery was not recognized for decades, and that it was ultimately validated by molecular methods that she did not use, is a testament to the power of her insights and the depth of her understanding of genetics.

McClintock's work also has a broader significance for the sociology of science. Her experience of isolation and incomprehension illustrates the difficulty that scientists can face when their discoveries challenge established paradigms, and it serves as a reminder that the scientific community is not always quick to recognize the importance of genuinely novel ideas. McClintock's perseverance in the face of skepticism---her refusal to abandon her work or to compromise her vision---is an inspiration to scientists everywhere.

The mobile genetic elements that McClintock discovered have turned out to be one of the major and ubiquitous features of genomes. They have shaped the evolution of life on Earth, they play important roles in human health and disease, and they have become powerful tools for genetic research and therapy. Barbara McClintock's discovery of transposable elements is, in the words of the Nobel Committee, ``one of the two or three most outstanding discoveries in the genetics of the twentieth century.''


\section{Modern Developments}

McClintock's transposon discovery has led to modern genomics and gene therapy. Transposons are now used as tools for gene delivery and insertional mutagenesis screens. Understanding that ~45% of the human genome consists of transposable elements has revealed their role in evolution and disease. LINE-1 retrotransposons are active in the human genome and may contribute to cancer and neurological diseases. CRISPR-Cas9 has been adapted from a bacterial immune system (related to transposon defense) for genome editing.


\section{Bibliography}

\begin{thebibliography}{9}
\bibitem{nobel-1983}
Nobel Prize Outreach, ``The Nobel Prize in Physiology or Medicine 1983,''\emph{NobelPrize.org}.\\
\url{https://www.nobelprize.org/prizes/medicine/1983/summary/}

\bibitem{lecture-1983}
Barbara McClintock, ``Nobel Lecture,''\emph{NobelPrize.org}. Winner-authored primary source; downloadable from the official Nobel Prize archive.\\
\url{https://www.nobelprize.org/prizes/medicine/1983/mcclintock/lecture/}
\end{thebibliography}
